\section{Common-shift defect one and residual families}
\label{sec:common-shift-residuals}

The following field-theoretic lemma isolates exactly what a common
translation can and cannot lose.  The independent anchor is essential: merely
excluding the possibility that the shift parameter belongs to the base field
would not prove the final two-variable independence statement.

\begin{theorem}[Common-shift defect one]
\label{thm:common-shift-defect}
Let \(K\subset\C\) be a field of characteristic zero, and let
\(z_0,z_1,\ldots,z_r\) be algebraically independent over \(K\), where
\(r\geq1\).  Fix \(c,e_0\in K^\times\), and put
\begin{equation}\label{eq:common-shift-coordinates}
 T_j=c(z_j-\gamma)\quad(1\leq j\leq r),
 \qquad
 \Delta=e_0(z_0-\gamma).
\end{equation}
Then
\begin{equation}\label{eq:common-shift-lower-bound}
 \trdeg_K K(T_1,\ldots,T_r)\geq r-1.
\end{equation}
More precisely, the \(r-1\) anchored differences
\[
 T_2-T_1,\ldots,T_r-T_1
\]
are algebraically independent over \(K\).  Consequently the \(T_j\) are
algebraically dependent over \(K\) if and only if
\[
 \trdeg_K K(T_1,\ldots,T_r)=r-1.
\]
In that case \(\gamma\) and \(\Delta\) are algebraically independent over
\(K\).
\end{theorem}

\begin{proof}
The linear forms \(z_j-z_1\), \(2\leq j\leq r\), have row rank \(r-1\)
in the independent coordinates \(z_1,\ldots,z_r\).  Hence
\[
 T_j-T_1=c(z_j-z_1),\qquad 2\leq j\leq r,
\]
are algebraically independent over \(K\), proving
\eqref{eq:common-shift-lower-bound}.

Suppose now that the \(T_j\) are algebraically dependent and write
\(L=K(T_1,\ldots,T_r)\).  The lower bound gives
\(\trdeg_K L=r-1\).  Since \(z_j=\gamma+c^{-1}T_j\), one has
\begin{equation}\label{eq:common-shift-field-identity}
 K(z_1,\ldots,z_r,\gamma)=L(\gamma).
\end{equation}
The left side contains the purely transcendental field
\(K(z_1,\ldots,z_r)\), of transcendence degree \(r\), whereas the right
side has transcendence degree at most \((r-1)+1=r\).  Thus both sides of
\eqref{eq:common-shift-field-identity} have transcendence degree exactly
\(r\).  It follows that \(\gamma\) is algebraic over
\(K(z_1,\ldots,z_r)\), and also that \(\gamma\) is transcendental over
\(L\), hence over \(K\).

Because \(z_0,z_1,\ldots,z_r\) are algebraically independent, \(z_0\) is
transcendental over the algebraic closure of
\(K(z_1,\ldots,z_r)\).  In particular it is transcendental over
\(K(\gamma)\).  Therefore \(z_0-\gamma\) is transcendental over
\(K(\gamma)\).  Since \(\gamma\) is transcendental over \(K\) and
\(e_0\in K^\times\), the two numbers
\(\gamma\) and \(\Delta=e_0(z_0-\gamma)\) are algebraically independent
over \(K\).
\end{proof}

We apply the theorem over \(K=\Eexp\).  For \(N,n\geq1\), define
\begin{align}
 Y_N&=(N+1)\Ein(2^N)-N\Ein(2^{N+1}),
 &\xi_n&=\Ein(-n),\label{eq:residual-Ein-coordinates}\\
 R_N&=e\bigl((N+1)E_1(2^N)-N E_1(2^{N+1})\bigr),
 &S_n&=e\bigl(\Ei(n)-\log n\bigr).
 \label{eq:residual-classical-coordinates}
\end{align}
The connection formulas give
\begin{equation}\label{eq:residual-common-shift}
 R_N=e(Y_N-\gamma),
 \qquad
 -S_n=e(\xi_n-\gamma),
 \qquad
 \delta=e(\Ein(1)-\gamma).
\end{equation}

\begin{lemma}[Independent residual coordinates]
\label{lem:residual-Ein-independent}
The countable family
\[
 \{\Ein(1)\}\cup\{Y_N:N\geq1\}\cup\{\xi_n:n\geq1\}
\]
is algebraically independent over \(\Eexp\), in the finite-subset sense.
\end{lemma}

\begin{proof}
The pure-exponential base-change theorem makes the distinct coordinates
\[
 \Ein(1),\qquad \Ein(2^k)\ (k\geq1),\qquad \Ein(-n)\ (n\geq1)
\]
algebraically independent over \(\Eexp\).  The negative coordinates are
disjoint from the positive ones.  It remains only to check the row rank of
the \(Y_N\).  In a nonzero finite linear combination of their defining
rows, choose the smallest occurring index \(N\).  The coordinate
\(\Ein(2^N)\) occurs only in that smallest row and has coefficient \(N+1\),
so its row coefficient must vanish.  Iteration proves that every row
coefficient vanishes.  Lemma~\ref{lem:linear-image-ein} now gives the
assertion.
\end{proof}

\begin{corollary}[Countable residual dichotomy]
\label{cor:countable-residual-dichotomy}
Exactly one of the following alternatives holds.
\begin{enumerate}[label=\textup{(\roman*)}]
 \item The family
 \[
  \{R_N:N\geq1\}\cup\{S_n:n\geq1\}
 \]
 is algebraically independent over \(\Eexp\).
 \item Some finite subfamily is algebraically dependent over \(\Eexp\),
 and \(\gamma,\delta\) are then algebraically independent over
 \(\Eexp\).  In particular \(e,\gamma,\delta\) are algebraically
 independent over \(\Qbar\).
\end{enumerate}
Unconditionally, the anchored family
\begin{equation}\label{eq:anchored-residual-family}
 \boxed{
 \{R_N-R_1:N\geq2\}\cup\{S_n+R_1:n\geq1\}}
\end{equation}
is algebraically independent over \(\Eexp\).
\end{corollary}

\begin{proof}
Apply Theorem~\ref{thm:common-shift-defect} to every finite selection from
\eqref{eq:residual-common-shift}, with
\(z_0=\Ein(1)\), common scale \(c=e\), and \(\Delta=\delta\).
Lemma~\ref{lem:residual-Ein-independent} verifies the hypothesis.  A
countable family is algebraically dependent exactly when one of its finite
subfamilies is dependent, which proves the dichotomy.

For the final assertion, anchor the common-shift variables at \(R_1\).  Their
differences are
\[
 R_N-R_1=e(Y_N-Y_1),
 \qquad
 (-S_n)-R_1=-(S_n+R_1),
\]
and the anchored-difference assertion of
Theorem~\ref{thm:common-shift-defect} proves their joint algebraic
independence.
\end{proof}

\begin{remark}[Logical scope]
The corollary is an exact alternative, not a construction of the relation in
its second branch.  If no finite residual relation is supplied, the result is
the unconditional algebraic independence of the anchored family
\eqref{eq:anchored-residual-family}; it does not decide \(\gamma\) or
\(\delta\) individually.
\end{remark}
